\documentclass{article}
\usepackage[utf8]{inputenc}
\usepackage[left=3cm, right=3cm, top=2cm, bottom=2cm]{geometry}
\usepackage[linesnumbered,boxed]{algorithm2e}
\usepackage[noend]{algpseudocode}
\usepackage{amsmath}
\usepackage{amssymb}
\usepackage{amsthm}
\usepackage{authblk}
\usepackage{comment}
\usepackage{mdframed}
\usepackage{graphicx}

\newtheorem{definition}{Definition}[section]

\begin{document}


\title{Homework 3}
\author{CS 2051, Spring 2022}
\affil{Georgia Tech\\~\\{\bf Due February 12}}
\date{}
\maketitle


\paragraph{Reminder:} You may collaborate with other students in this
class, but (1) you must write up your own solutions in your own words, and (2) 
you must write down everyone you worked with at the top of the page, or ``no 
collaborators'' if you did it all on your own. Additionally, if you used any outside websites or textbooks besides the course text, please cite them here. You may not use question-answer sites like Chegg or MathOverflow. \\

\hline

\subsection*{Question 1}
Let $A, B \subseteq U$ where $U$ is the universe. 
\begin{enumerate}
    \item Prove that $(A \cup B)^C = A^C \cap B^c$.
        \begin{mdframed}
                        $(A \cup B)^c$ can be written as $\{x | x \in U \wedge x \notin A \cup B \}$ by the definition of complement:
                \begin{equation}
                    \{x | x \in U \wedge x \notin A \cup B\}
                \end{equation}
                \begin{equation}
                    \{x | x \in U \wedge \neg (x \in A \cup B)\}\text{, by definition of negation of sets}
                \end{equation}
                \begin{equation}
                    \{x | x \in U \wedge \neg (x \in A \vee x \in B)\}\text{, by definition of union}
                \end{equation}
                \begin{equation}
                    \{x | x \in U \wedge \neg (x \in A) \wedge \neg (x \in B)\}\text{, by DeMorgan's law for logical equivalences}
                \end{equation}
                \begin{equation}
                    \{x | x \in U \wedge x \notin A \wedge x \notin B\}\text{, by definition of negation}
                \end{equation}
                \begin{equation}
                    \{x | x \in U \wedge x \in A^c \wedge x \in B^c\}\text{, by definition of complement}
                \end{equation}
                \begin{equation}
                    \{x | x \in U \wedge x \in (A^c \cap B^c)\}\text{, by definition of intersection}
                \end{equation}
                Thus, $(A \cup B)^c=A^c \cap B^c$.
        \end{mdframed}

    \item Prove that $(A \cap B)^C = A^C \cup B^C$.
            \begin{mdframed}
                $(A \cup B)^c$ can be written as $\{x | x \in U \wedge x \notin A \cap B$ by the definition of complement:
                \begin{equation}
                    \{x | x \in U \wedge x \notin A \cap B\}
                \end{equation}
                \begin{equation}
                    \{x | x \in U \wedge \neg (x \in A \cap B)\}\text{, by definition of negation of sets}
                \end{equation}
                \begin{equation}
                    \{x | x \in U \wedge \neg (x \in A \wedge x \in B)\}\text{, by definition of intersection}
                \end{equation}
                \begin{equation}
                    \{x | x \in U \wedge \neg (x \in A) \vee \neg (x \in B)\}\text{, by DeMorgan's law for logical equivalences}
                \end{equation}
                \begin{equation}
                    \{x | x \in U \wedge x \notin A \vee x \notin B\}\text{, by definition of negation}
                \end{equation}
                \begin{equation}
                    \{x | x \in U \wedge x \in A^c \vee x \in B^c\}\text{, by definition of complement}
                \end{equation}
                \begin{equation}
                    \{x | x \in U \wedge x \in (A^c \cup B^c)\}\text{, by definition of union}
                \end{equation}
                Thus, $(A \cup B)^c=A^c \cap B^c$.
            \end{mdframed}
    
    \item Prove that
    \begin{gather*}
            (X \cap Y \cap Z) \cup (X \cup Y \cup Z)^C \\
            \text{and} \\
            (X \cup Y^C) \cap (X \cup Z^C) \cap (Y \cup X^C) \cap (Y \cup Z^C) \cap (Z \cup X^C) \cap (Z \cup Y^C) \\
            \text{are equal.}
    \end{gather*}
            \begin{mdframed}
            We can rewrite the LHS (left-hand side) as: \\
            \begin{gather*}
            (\{ x | x \in X \} \wedge \{ x | x \in Y \} \wedge \{ x | x \in Z \}) \vee \neg ((\{ x | x \in X \} \vee \{ x | x \in Y \} \vee \{ x | x \in Z \})
            \end{gather*} using its membership function.
            Then we can DeMorgan it:
            \begin{gather*}
                (\{ x | x \in X \} \wedge \{ x | x \in Y \} \wedge \{ x | x \in Z \}) \vee ((\neg\{ x | x \in X \} \wedge \neg\{ x | x \in Y \} \wedge \neg\{ x | x \in Z \})
            \end{gather*}
            And then we can distribute it:
            \begin{gather*}
                (\{ x | x \in X \} \vee \neg\{ x | x \in X \}) \wedge (\{ x | x \in X \} \vee \neg\{ x | x \in Y \}) \wedge (\{ x | x \in X \} \vee \neg\{ x | x \in Z \}) \wedge \\
                (\{ x | x \in Y \} \vee \neg\{ x | x \in X \}) \wedge (\{ x | x \in Y \} \vee \neg\{ x | x \in Y \}) \wedge (\{ x | x \in Y \} \vee \neg\{ x | x \in Z \}) \wedge \\
                (\{ x | x \in Z \} \vee \neg\{ x | x \in X \}) \wedge (\{ Z | x \in X \} \vee \neg\{ x | x \in Y \}) \wedge (\{ x | x \in Z \} \vee \neg\{ x | x \in Z \})
            \end{gather*}
            By the law of the excluded middle ($p \vee \neg p$ is true), we can reduce this to the RHS by definition of the membership function:
            \begin{gather*}
                (\{ x | x \in X \} \vee \neg\{ x | x \in Y \}) \wedge (\{ x | x \in X \} \vee \neg\{ x | x \in Z \}) \wedge \\
                (\{ x | x \in Y \} \vee \neg\{ x | x \in X \}) \wedge (\{ x | x \in Y \} \vee \neg\{ x | x \in Z \}) \wedge \\
                (\{ x | x \in Z \} \vee \neg\{ x | x \in X \}) \wedge (\{ Z | x \in X \} \vee \neg\{ x | x \in Y \}) \\
                = (X \cup Y^C) \cap (X \cup Z^C) \cap (Y \cup X^C) \cap (Y \cup Z^C) \cap (Z \cup X^C) \cap (Z \cup Y^C)
            \end{gather*}
            \end{mdframed}

\end{enumerate}


\subsection*{Question 2}
For each of the following parts, determine if each of the following are true: $A \subseteq B$, $B \subseteq A$, $A = B$. Prove your answers. 
\begin{enumerate}
    \item $A = \{2n : n \in \mathbb{Z} \}, B = \{4n : n \in \mathbb{Z} \}$
    \begin{mdframed}
        \begin{enumerate}
            \item $A \not\subseteq B$: Consider 2. $2 \in A$ because $2(1) = 2$, but $2 = 4(1/2)$ and $1/2 \notin \mathbb{Z}$, so $2 \notin B$.
            \item $B \subseteq A$: Let $y \in B$. so $y = 4k$ for some $k \in \mathbb{Z}$. $4k = 2(2k) = 2p$ for some $p \in \mathbb{Z}$. $2p \in A$ by definition, so $B \subset A$.
        \end{enumerate}
    \end{mdframed}
    \item $A = \mathbb{Q}, B = \{7n : n \in \mathbb{Q}\}$
    \begin{mdframed}
        \begin{enumerate}
            \item $A \subseteq B$: Consider $y \in A$. $y = \frac{p}{q}$ where $p, q \in \mathbb{Z}, q \neq 0$. Then multiply by $\frac{7}{7}$ to get $\frac{7p}{7q} \in A$. Let $a = 7q$, then $\frac{7p}{a}$, so $y \in B$ by definition.
            \item $B \subseteq A$: Consider $x \in B$. Then $x = 7n, n \in \mathbb{Q}$. Therefore, $x = 7\frac{p}{q}, p, q\in \mathbb{Z}, q \neq 0$ and $x = \frac{a}{q}$ if $a = 7p$. $a, q \in \mathbb{Z}$ and $q \neq 0$, so $x \in A$.
            \item $A = B$. Because $A \subseteq B$ and $B \subseteq A$, $A = B$.
        \end{enumerate}
    \end{mdframed}

    \item $A = \{2n : n \in \mathbb{Z}\}$, $B = \{n = 2k + 1 \text{ for some k in }\mathbb{Z} : n \in \mathbb{R} \}^C$
    \begin{mdframed}
        \begin{enumerate}
            \item $A \subseteq B$: Let $y \in A$. Then $y = 2k$ for $k \in \mathbb{Z}$. Assume towards contradiction that $y \notin B$, then there exists some $j \in \mathbb{Z}$ such that $2j + 1 = y$. This would mean that $2j + 1 = 2k$, which contradicts that $k$ is an integer becuase $k = \frac{2j + 1}{2} = j + \frac{1}{2}$. Therefore, $y \in B$ and so $A \subseteq B$.
            \item $B \not\subseteq A$: Observe that the universe of $B$ is $\mathbb{R}$, so the complement of "odd integers" is "all real numbers that are not odd integers". Therefore, $\pi \in B$. Because $\pi$ is irrational, there does not exist $k \in \mathbb{Z}$ such that $\pi = 2k$ becuase $\frac{\pi}{2} = k$ is not an integer.
        \end{enumerate}
    \end{mdframed}
\end{enumerate}

\subsection*{Question 3}
Out of the following four functions, one is a surjection only, one is an injection only, one is neither, and one is a bijection. State which is which and demonstrate your answers by (1) proving
that it is an injection or naming two distinct elements of the domain that are associated to the
same element of the target, and (2) proving that it is a surjection or naming an element of the
target for which there is no corresponding element in the domain.
\begin{enumerate}
    \item $f: \mathbb{Z}^2 \mapsto \mathbb{Z}$, $f(m, n) = m^2 + n^2$
    \begin{mdframed}
    This is neither.
    \begin{enumerate}
        \item Not an injection: $f(0, 1) = f(1, 0) = 1$.
        \item Not a surjection: Because the range of both $m^2$ and $n^2$ are $[0, +\infty)$, $f(m, n) \geq 0$ for all $m, n$. Therefore, the codomain is not equal to the range. For example, $-1$ is not in the codomain.
    \end{enumerate}
    \end{mdframed}
    \item $f: \mathbb{Z}^2 \mapsto \mathbb{Z}$, $f(m, n) = m$
    \begin{mdframed}
    This is a surjection only.
    \begin{enumerate}
        \item Not an injection: $f(0, 1) = f(0, 2) = 0$.
        \item Surjection: For any $z \in \mathbb{Z}$, we have that $f(z, 0) = z$. Therefore, every element in the codomain has at least one element in the domain so it is a surjection by definition.
    \end{enumerate}
    \end{mdframed}
    \item $f: \mathbb{R}^2 \mapsto \mathbb{C}$, $f(m, n) = m - 2ni$
    \begin{mdframed}
    This is a bijection.
    \begin{enumerate}
        \item Injection: Let $(m, n), (x, y) \in \mathbb{R}^2$ where $f(m, n) = f(x, y)$. For $f$ to be an injection we want to show that, $(m, n) = (x, y)$. \\
        $m = 2ni = x - 2yi$, so $m = x$ and $-2ni = -2yi \rightarrow ni = yi \rightarrow n = y$. This is because real numbers are closed under addition and multiplication by a real number and imaginary numbers are closed under addition and multiplication by a real number.
        \item Surjection: Let $a + bi \in \mathbb{C}$ by any member of the range. Then consider $(a, \mathbb{-b}{2} \in \mathbb{R^2}$. Then $f(a, -b/2) = a + bi$, so there exists an element in the domain that maps to $a + bi$ for any $a + bi \in \mathbb{C}$.
    \end{enumerate}
    \end{mdframed}
    \item $f: \mathbb{Z}^2 \mapsto \mathbb{Z}$, $f(m, n) = 2^m3^n$
    \begin{mdframed}
    \begin{enumerate}
        \item Injection: Let $(m, n), (x, y) \in \mathbb{R}^2$ where $f(m, n) = f(x, y)$. Then $2^m3^n = 2^x3^y$ and therefore, $2^{m-x} = 3^{y-n}$. This is possible iff $m - x = y - n = 0$, so $m = x$ and $y = n$. Therefore, $f$ is an injection.
        \item Not a surjection: $2^m3^n$ is always strictly positive, so there is no $(m, n)$ such that $f(m, n) = 0$.
    \end{enumerate}
    \end{mdframed}

\end{enumerate}


\subsection*{Question 4}
Determine whether the relation $\mathcal{R}$ on the set of all people is reflexive, symmetric, antisymmetric, and/or transitive. You do not need to justify your answers.\\

$(a,b)\in \mathcal{R}$ if and only if
\begin{enumerate}
\item $a$ is taller than $b$.\\
$\bigcirc$ \texttt{Reflexive} \quad $\bigcirc$ \texttt{Symmetric} \quad $\bigcirc$ \textcolor{blue}{\texttt{Antisymmetric} \quad $\bigcirc$ \texttt{Transitive}}

\item $a$ and $b$ were born on the same day.\\
\textcolor{blue}{$\bigcirc$ \texttt{Reflexive} \quad $\bigcirc$ \texttt{Symmetric}} \quad $\bigcirc$ \texttt{Antisymmetric} \quad $\bigcirc$ \textcolor{blue}{\texttt{Transitive}}

\item $a$ has the same first name as $b$.\\
\textcolor{blue}{$\bigcirc$ \texttt{Reflexive} \quad $\bigcirc$ \texttt{Symmetric}} \quad $\bigcirc$ \texttt{Antisymmetric} \quad $\bigcirc$ \textcolor{blue}{\texttt{Transitive}}

\item $a$ and $b$ have a common grandparent.\\
\textcolor{blue}{$\bigcirc$ \texttt{Reflexive} \quad $\bigcirc$ \texttt{Symmetric}} \quad $\bigcirc$ \texttt{Antisymmetric} \quad $\bigcirc$ \texttt{Transitive}

\item $a$ and $b$ have met.\\ 
$\bigcirc$ \texttt{Reflexive} \quad $\bigcirc$ \textcolor{blue}{\texttt{Symmetric}} \quad $\bigcirc$ \texttt{Antisymmetric} \quad $\bigcirc$ \texttt{Transitive}

\end{enumerate}
\textbf{Note: for \#5, both "reflexive" and "not reflexive" were accepted.}

\subsection*{Question 5}
The set of rational numbers can be seen as a set of the equivalence classes of a certain relation on $\mathbb{Z} \times \mathbb{Z}$ defined by $(a, b)R(c, d)$ if and only if $ad = bc$. Show that $R$ is an equivalence relation.

\begin{mdframed}
\begin{enumerate}
    \item \textbf{Reflexive.} Clearly, $ab=ba$ since multiplication is commutative.
    \item \textbf{Symmetric.} $(a,b)\mathcal{R} (c,d)$ iff $ad=bc$ and then $cb=da$ which is equivalent to $(c,d)\mathcal{R} (a,b)$.
    \item \textbf{Transitive.} Let $(a,b)\mathcal{R} (c,d)$ and $(c,d)\mathcal{R} (e,f)$. Then $ad=bc$ and $cf=de$ or, $\frac{a}{b}=\frac{c}{d}=\frac{e}{f}$ yielding $af=be$ or  $(a,b)\mathcal{R} (e,f)$.
\end{enumerate} \\

\textbf{Note: Technically the transitive case only holds when $b, d \neq 0$. No points were deducted for not observing this.}
\end{mdframed}



\subsection*{Question 6 \textit{(Hilbert's Grand Hotel)}}
     The Grand Hotel has countable many rooms, numbered 1, 2, 3, \dots, and it is fully occupied.
    \begin{enumerate}
        \item Suppose that $k$ guests arrive at the hotel. Show how you can accommodate all new guests without removing any of the current guests.
    \begin{mdframed}
    Shift all existing guests $k$ rooms over, and place the $k$ guests in the rooms $1, \dots k$.
    \end{mdframed}
        \item Suppose that a countable infinite number of guests arrive. Show how you can accommodate all new guests without removing any of the current guests.
    \begin{mdframed}
    Move all existing guests in room $n$ to room $2n + 1$. This frees all rooms that can be expressed as $2n$ for the countably infinite number of guests to occupy.
    \end{mdframed}
        \item Suppose Hilbert (the owner of the Grand Hotel) expands the property to a second building, also with infinite rooms numbered $1, 2, 3, \dots$. Show how you can spread the current guests to fill every room in both buildings.
    \begin{mdframed}
    Take every guest with an odd room number $2k - 1$, where $k \in \mathbb{N}$ and move them into room number $k$ in the second building. \\
        Now take the remaining guests in the original building. Since their room numbers are even, we can represent their room number as $m = 2l$, where $l \in \mathbb{N}$. Move these guests into room number $l$ in the first building.
    \end{mdframed}
    \end{enumerate}

\subsection*{Question 7 (\textit{Cantor's theorem})}



\begin{enumerate}
    \item Prove that there does not exist an onto function between $\mathbb{N}$ and its power set, $\mathcal{P}(\mathbb{N})$.\\

 {\it Hint: proceed by contradiction assuming there is such function $f:\mathbb{N} \rightarrow \mathcal{P}(\mathbb{N})$. Let 
 
 \[
 T=\{n\in \mathbb{N}|~n\notin f(n) \}.
 \]

Since $f$ is onto, there exist $t\in \mathbb{N}$ such that $f(t)=T$. Argue a contradiction by looking at $t\in T$ and $t\notin T$.}

\begin{mdframed}
Assume towards contradiction that there exists a surjection $f: \mathbb{N} \mapsto \mathbb{P}(\mathbb{N})$. Let $T = \{n \in \mathbb{N} | n \notin f(n) \}$ $T \in \mathbb{P}(\mathbb{N})$, so there must exist $t \in \mathbb{N}$ such that $f(t) = T$ becuase $f$ is a surjection. \\
Now, consider two cases, where $t \in f(t)$ and where $t \notin f(t)$:
\begin{enumerate}
    \item $t \in T$: Then $t \notin f(t)$, but $f(t) = T$ by definition. But $t \in \mathbb{P}(\mathbb{N})$ becuase the power set contains all possible subsets of $\mathbb{N}$, so $t \in f(t)$, which is a contradiction.
    \item $t \notin T$: Then $t \in f(t)$, but because $f(t) = T$, $t \in T$, which is a contradiction.
\end{enumerate} 
These cover all possible cases of $t$, so we have a contradiction and therefore there is no onto function.
\end{mdframed}


\item Show that the set of all countable sets is uncountable

\begin{mdframed}
$\mathbb{P}(\mathbb{N})$ is a subset of the set of all countable sets. $\mathbb{P}(\mathbb{N})$ is uncountable, so any supersets of $\mathbb{P}(\mathbb{N})$ must also be uncountable. Therefore, the set of all countable sets is uncountable.
\end{mdframed}


\end{enumerate}



\end{document}